\documentclass[pdftex,12pt,letter]{article}
\usepackage[pdftex]{graphicx}

\usepackage{pdfpages} 
\usepackage{fixltx2e}

\addtolength{\textwidth}{3.5cm}
\addtolength{\hoffset}{-2cm}
\addtolength{\textheight}{4.5cm}
\addtolength{\voffset}{-3cm}

\bibliographystyle{plain}
\usepackage{natbib}
\newcommand{\HRule}{\rule{\linewidth}{0.5mm}}

\renewcommand{\thesection}{\Alph{section}}

%\usepackage{blindtext}
\usepackage{listings}

\usepackage{amsmath}

\usepackage{subfig}
\begin{document}
\begin{titlepage}

\begin{center}


% Upper part of the page
\includegraphics[width=.5\textwidth]{./logo}\\[1cm]    

\textsc{\LARGE CS-GY 6903: Modern Cryptography}\\[1.5cm]

\textsc{\Large Professor: Giovanni Di Crescenzo}\\[0.5cm]

% Title
\HRule \\[0.4cm]
{ \large \bfseries Project 1:\\
{\small Breaking Polyalphabetic ciphers}}\\[0.4cm]

\HRule \\[1.5cm]
% first column
\begin{minipage}{0.4\textwidth}
\begin{flushleft}
\large Santiago Torres-Arias\\
{\small {\bf NYU ID: N14553751 } }
\end{flushleft}
\end{minipage}
%second column
\begin{minipage}{0.40\textwidth}
\begin{flushright}
\large Lucas Mladek\\
{\small {\bf NYU ID: N1xxxxxx } } 
\end{flushright}
\end{minipage} 

\vfill

% Bottom of the page
{\large \today}

\end{center}
\end{titlepage}

\newpage
\tableofcontents
\newpage

\section{Introduction}

Polyalphabetic ciphers have been around in humanity for hundreds of years now.
And, although their use is no longer recommended for real cryptography, its use
is still widespread among cryptography amateurs and cryptography
challenges\footnote{A really popular cipher challenge includes two
polyalphabetic ciphers http://simonsingh.net/cryptography/cipher-challenge/}.

It is known that polyalphabetic ciphers are still subject to frequency
analysis, and known plaintext attacks. The focus of this writeup is to
introduce techniques that leverage this knowledge to break these ciphers in the
least possible time.

\section{Understanding polyalphabetic ciphers}
% Here, we talk about what's a polyalphabetic cipher
% and how the guy must be high about thinking of the j(x) function


\subsection{The anatomy of the $j(i)$ function}

We can try to decompose the $j(i)$ function by using calculus concepts. The
following equation provides a simplified version of the possible variants that
$j(i)$ can have:

\begin{equation}
    j(i) = \left( \frac{xi^y}{zj^k} + m \right) \mod{t}
\end{equation}

In this equation, we adopt the following nomenclature:
\begin{itemize}
    \item i is the index
    \item x is a define a ``decimating coefficient'' for i.
    \item y defines a ``decimating exponential'' factor for the index.
    \item function $zj^k$ defines a ``stretching factor''
    \item m is a ``start offset'', and it can be described as shifting the key by m bins.
    \item $\mod{t}$ guarantees that the values selected are inside the generated key and t is the length of the key.
\end{itemize}

The relationship between the upper and lower factors define how this function behaves. 
In general, we can consider to be three variants: periodic, decimating and stretching. 
We will define these concepts next.

\subsection{Key stretching and key-decimating}

When the factor above ($xi^y$) is polinomially bigger than the factor below
($zj^k$), then we can assume that we have a key-decimating function. Key
decimating functions make the key essentially "smaller" in a sense that some
values are skipped and hence lost (depending on the periodicity, we will
explain this later).

In the other hand, when values from the upper monomial are smaller than the
lower one, we have what's called a key-stretching function. When this happens,
we see that -- due to integer math -- some values are repeated contiguously.
For example, imagine that we have the following $j(i)$:

\begin{equation}
    j(i) = \frac{i}{4}\mod{30}
\end{equation}

In this case, the resulting function can be considered stretching because
values for i below 4 will all point to 0, which is the first element of the
key.

The nature of this equation has a direct effect on what we will call the key's
periodicity. We will explain this in the next subsection.

\subsection{Periodicity}

Periodicity of the key (and $j(i)$) can be understood as ``the number of values for i
before the key-selecting sequence repeats''. Understanding this can allow us for 
an easy definition of a break function. 

Regular periodic functions can be defined as those functions in which the lower
factor is a multiple of t. When this happens, the periodicity of $j(i)$ also falls
in a multiple of t. For our breaking scheme, we considered the function $j(i)$ to  
be a regular periodic function since it was the easiest to analyze. However, in 
Appendix B, we describe methods that could have been used to tackle the non-regular
periodic family of functions.

\subsection{Why periodicity matters and $j(i)$ not as much}

If we consider a function to be regular periodic, then we can assume that the 
set of keys ${a...z}^t$ under a specific $j(i)$ is only a different permutation
of another set of keys under a different $j(i)$. Knowing this, we can
assume that, when a function is regular periodic, obtaining a proposed key for
$j(i) = i \mod(t)$ is a safe bet. We will use this assumption in our breaking 
mechanisms, as it simplifies our analysis greatly.

\section{Breaking dictionary 1}

In order to design the best-fastest mechanism to break the defined cipher, we 
first need to understand the nature of the existing dictionaries. In this case
we know that the existing plaintexts falls in a range of 150 different plaintexts
so we decided to analyse them as words.

\subsection{Minimum prefix: plaintexts as words}

Since the set of possible messages from Dictionary one is the same as the number of 
entries a dictionary 1, cryptanalysis is really easy. We could consider that the 
messages from dictionary one are a single, long word encrypted and sent through the wire.

To understand this, we build a python script called "minimum\_prefix.py", that analyses the 
words in a dictionary to identify the minimum amount of characters needed to differentiate
ciphertext from each other. The result of this script threw a value of 9, and we used it
as part of our dictionary-header definition; this script can be found in the sources
and is also described in Appendix A.

\subsection{Defining our oracle}

We can obtain a possible key by substracting a piece of plaintext from the ciphertext:

\begin{equation}
    p_{[i]} + k_{[x]} = c_{[i]} \rightarrow k_{[x]} = c_{[i]} - p_{[i]}
\end{equation}

Now, how do we know our key is correct?

\subsection{Breaking dictionary 1: our attack}

Breaking dictionary 1 is easy: we only need to apply the prefix plaintext to
the provided ciphertext and compare the resulting  key to another portion of
the plaintext. The only catch is that we need to identify a region that uses
the same range of the key. Since we assumed $j(i)$ to be regular periodic, then
we can assume that it is a place where period starts and verify if the result
matches correct plaintext of the same dictionary entry in that location. 

%should we talk about the speedups? this is getting complex enough by itself...
The code we wrote for this takes less than a second to verify and dismiss all
possible candidates for the plaintext. 



\section{Breaking dictionary 2}

\subsection{Case study: Rosignol Cipher - The man with the Iron mask}
https://en.wikipedia.org/wiki/Great\_Cipher

\subsection{Case Study: Breaking Indian Ciphers}
http://shodhganga.inflibnet.ac.in/bitstream/10603/2193/13/13\_chapter\%205.pdf

\subsection{Breaking polyalphabetic ciphers using triads(or trillables)}

\subsubsection{The size of the set}

\subsubsection{Probabilistic results}

\section{Results}

\section{Conclusion}

\newpage
\appendix{}
\section*{Appendix A: Statistical analysis scripts and code}
\end{document}
